% !PRINTED 75
\weberpage{75}
%<<CONT
ist. Um die Bedingungen f\"ur diesen Fall zu ermitteln, erweitern
wir den Begriff der \gesperrt{Resolventen}, wie wir ihn schon im \S.\,17
betrachtet haben, noch etwas.

Wir lassen $n$ die Reihe der Zahlen der Gruppe $\mathfrak{N}$ durchlaufen
 und bezeichnen mit $\chi(n)$ einen der Charaktere dieser
Gruppe. Die erweiterten Resolventen sind dann, wenn $r$ eine
$m^{\text{te}}$ Einheitswurzel ist:
\weq{5}{(\chi, r) = \sum^{n} \chi(n) r^n.}
Verstehen wir unter $\mathfrak{B}$ die zu $\mathfrak{A}$ reciproke Gruppe und lassen
$b$ die Zahlen von $\mathfrak{B}$ durchlaufen, so giebt es nach \S.\,12, 7. eine
dem Grade von $\mathfrak{B}$ gleiche Anzahl von Charakteren $\chi_b$, die dadurch
 ausgezeichnet sind, dass
\[
\chi_b(a) = 1
\]
ist, f\"ur jede Zahl $a$ aus der Gruppe $\mathfrak{A}$. Daraus folgt dann nach
der Definition der Charaktere \S.\,11, (4) f\"ur jedes $n$:
\weq{6}{\chi_b(a n) = \chi_b(n),}
d. h. der Charakter $\chi_b$ hat f\"ur alle Zahlen einer jeden Nebengruppe
 $\mathfrak{A} n_1, \mathfrak{A} n_2 \ldots$ einen und denselben Werth. Demnach wird
die Resolvente $(\chi_b, r)$ nur von den Perioden $\eta$ abh\"angen, und
den Ausdruck erhalten:
\weq{7}{(\chi_b, r) = \eta + \chi_b(n_1) \eta_1 + \cdots + \chi_b(n_{e-1}) \eta_{e-1}.}
Wenn $\eta$ zu einer Gruppe $\mathfrak{A}'$ geh\"ort, so geh\"oren die conjugirten
 Zahlen $\eta_1, \eta_2, \ldots \eta_{e-1}$ zu derselben Gruppe (weil $\mathfrak{A}'$ ein
Normaltheiler von $\mathfrak{N}$ ist, Bd.\ I, \S.\,154). Wenn also $a'$ irgend
eine Zahl aus $\mathfrak{A}'$ ist, so bleiben die Gr\"ossen $\eta, \eta_1, \eta_2, \ldots \eta_{e-1}$
durch die Substitution $(r, r^{a'})$ unge\"andert und nach (7) ist auch
\weq{8}{(\chi_b, r) = (\chi_b, r^{a'}).}
Andererseits erh\"alt man aus (5), wenn man bedenkt, dass
$a' n$ zugleich mit $n$ die ganze Gruppe $\mathfrak{N}$ durchl\"auft,
\[
\begin{aligned}
(\chi_b, r) &= \sum^{n} \chi_b(n a') r^{n a'} = \chi_b(a') \sum^{n} \chi_b(n) r^{n a'}\\
&= \chi_b(a') (\chi_b, r^{a'}),
\end{aligned}
\]
also nach (8):
\weq{9}{(\chi_b, r) = \chi_b(a') (\chi_b, r).}
Ist $\mathfrak{A}$ nicht mit $\mathfrak{A}'$ identisch, sondern ein echter Theiler
von $\mathfrak{A}'$, so ist die zu $\mathfrak{A}'$ reciproke Gruppe $\mathfrak{B}'$ ein echter Theiler
von $\mathfrak{B}$ (nach \S.\,12, 8.); wenn also $b$ eine Zahl in $\mathfrak{B}$ ist, die nicht
%CONT>>
